\textbf{Proof.}
Work on the probability-one event on which thm:plugin-correspondence holds, and let \(U\) be any neighborhood of \(B^\star\).  The asserted isolation gives
\[
\gamma_U:=\inf_{B\in\mathcal B_\kappa\setminus U}w(B)-v_{\mathrm{mm}}>0.
\tag{1}
\]
Let \(\delta_m=\sup_B|\widehat w(B)-w(B)|\), so \(\delta_m\to0\).  If \(2\delta_m<\gamma_U\), then every \(B\notin U\) satisfies
\[
\widehat w(B)\ge w(B)-\delta_m
\ge v_{\mathrm{mm}}+\gamma_U-\delta_m
>v_{\mathrm{mm}}+\delta_m
\ge\widehat w(B^\star).
\]
Consequently \(\widehat{\mathcal A}_{\mathrm{mm}}\subset U\) eventually.  Since \(U\) was arbitrary, every measurable selection converges almost surely to \(B^\star\).  The unique active endpoint labels the limiting saddle; once singleton isolation is assumed, it is not needed for convergence of the basis itself.

For the failure of point convergence under ties, let \(n=3\), take \(f_i(w)=w_i\), and let the coordinates of \(W\) be independent Bernoulli \(1/2\).  Then
\[
\Pi=\frac12I_6,
\qquad
\Omega=\frac14\begin{bmatrix}I_3&-I_3\\-I_3&I_3\end{bmatrix}.
\tag{2}
\]
Let \(F_1=F_0=\delta_0\), let \(H=\delta_{(0,0)}\), and use the corresponding unpaired pilots.  Then
\[
\mu_1=\mu_0=\sigma_1^2=\sigma_0^2=c_L=c_U=0,
\qquad Q(0)=0.
\tag{3}
\]
Choose \(\kappa=1/8\) and, with \(e_j\) the coordinate vectors of \(\mathbb R^3\), define
\[
B^{(2)}=B(\sqrt3e_1,\sqrt3e_2),
\qquad B^{(3)}=B(\sqrt3e_1,\sqrt3e_3).
\]
For \(j\in\{2,3\}\), (2) gives
\[
S_{B^{(j)}}=\frac34I_2,
\qquad
\lambda_{\min}\{3^{-1}S_{B^{(j)}}\}=\frac14>\kappa.
\tag{4}
\]
Thus the two displayed bases are distinct feasible projective points.  By (3), however, \(h(0)=g_B(0)=\ell_B(0)=w(B)=0\) for every feasible \(B\).  Every pilot observation is identically zero, so the same equalities hold empirically and
\[
\widehat{\mathcal A}_{\mathrm{mm}}
=\mathcal A_{\mathrm{mm}}=\mathcal B_\kappa
\quad\text{at every pilot size}.
\tag{5}
\]
Along \(m_1=m_0=m\), select \(B^{(2)}\) for even \(m\) and \(B^{(3)}\) for odd \(m\).  This deterministic measurable selection has two distinct subsequential limits, while (5) makes the outer-set conclusion exact. \(\square\)